% GENERATED FILE. Edit canonical lessons, then run npm run generate:books.
% canonical-source-hash: fnv1a64:33ef90d775c84945
% canonical-lessons: IT-C428-decollare, IT-C428-atterrare, IT-C428-turbolenza, IT-C428-allacciare, IT-C428-finestrino, IT-R428-ripasso-421-428

\chapter{In volo}
\label{ch:it-in-volo}
\hlchaptermodality{428}

\begin{chapteropening}
\textbf{By the end of this chapter:} \emph{I can talk about a flight: taking off, landing, turbulence, fastening my belt and the window seat.}

Complete the last lesson of chapter 428: I can talk about a flight: taking off, landing, turbulence, fastening my belt and the window seat.
\end{chapteropening}

\section[decollare]{decollare — to take off (of a plane)}

\label{lesson:IT-C428-decollare}



\begin{quote}
Before the new one: say the Italian for the search, then the Italian for the pass.
\end{quote}

\subsection*{You'll want to know: decollare}
\textbf{decollare} — \textquotedblleft{}to take off (of a plane)\textquotedblright{}.

\begin{quote}
L'aereo decolla tra cinque minuti.
\end{quote}

\begin{quote}
\emph{The plane takes off in five minutes.}
\end{quote}

\emph{Il decollo} is take-off; figuratively, \emph{il progetto è decollato} — the project took off.

\begin{cousinweb}
\emph{decollare} was borrowed from French \emph{décoller} 'to come unstuck', built on \emph{colle} 'glue', from Greek \emph{kólla}: the wheels come unstuck from the ground.
\end{cousinweb}

\begin{grammarlens}[title={What you've built}]
One more word for telling, explaining and describing at B1.
\end{grammarlens}

\subsection*{Your turn}
\begin{itemize}
\raggedright
  \item \emph{Say it:} \emph{decollare}
  \item \emph{Say it:} \emph{decollare}, once more
  \item \emph{Say it:} a sentence of your own with \emph{decollare}
  \item \emph{Recall:} say the Italian for the search, then the Italian for the pass, then say \emph{decollare} again
\end{itemize}

\begin{tcolorbox}[breakable,colback=teal!4,colframe=teal!35!black,title={Before you move on}]
What does decollare mean? (\textquotedblleft{}To take off (of a plane)\textquotedblright{}.) Say it once more.
\end{tcolorbox}
\section[atterrare]{atterrare — to land (of a plane)}

\label{lesson:IT-C428-atterrare}



\begin{quote}
Before the new one: say the Italian for the pass, then the Italian for to take off.
\end{quote}

\subsection*{You'll want to know: atterrare}
\textbf{atterrare} — \textquotedblleft{}to land (of a plane)\textquotedblright{}.

\begin{quote}
Siamo atterrati in orario.
\end{quote}

\begin{quote}
\emph{We landed on time.}
\end{quote}

\emph{L'atterraggio} — the landing. Built on \emph{terra}.

\begin{cousinweb}
\emph{atterrare} is built from \emph{a-} 'to' and \emph{terra} 'ground', which is inherited from Latin \emph{terra} (EtymDB).
\end{cousinweb}

\begin{grammarlens}[title={What you've built}]
One more word for telling, explaining and describing at B1.
\end{grammarlens}

\subsection*{Your turn}
\begin{itemize}
\raggedright
  \item \emph{Say it:} \emph{atterrare}
  \item \emph{Say it:} \emph{atterrare}, once more
  \item \emph{Say it:} a sentence of your own with \emph{atterrare}
  \item \emph{Recall:} say the Italian for the pass, then the Italian for to take off, then say \emph{atterrare} again
\end{itemize}

\begin{tcolorbox}[breakable,colback=teal!4,colframe=teal!35!black,title={Before you move on}]
What does atterrare mean? (\textquotedblleft{}To land (of a plane)\textquotedblright{}.) Say it once more.
\end{tcolorbox}
\section[la turbolenza]{la turbolenza — the turbulence}

\label{lesson:IT-C428-turbolenza}



\begin{quote}
Before the new one: say the Italian for to take off, then the Italian for to land.
\end{quote}

\subsection*{You'll want to know: la turbolenza}
\textbf{la turbolenza} — \textquotedblleft{}the turbulence\textquotedblright{}.

\begin{quote}
Durante il volo c'è stata un po' di turbolenza.
\end{quote}

\begin{quote}
\emph{During the flight there was a bit of turbulence.}
\end{quote}

\emph{Turbolento} also describes a restless person.

\begin{cousinweb}
\emph{turbolenza} is a learned borrowing from Latin \emph{turbulentia}, from \emph{turbulentus} 'restless', built on \emph{turba} 'commotion, crowd' (Lewis \& Short).
\end{cousinweb}

\begin{grammarlens}[title={What you've built}]
One more word for telling, explaining and describing at B1.
\end{grammarlens}

\subsection*{Your turn}
\begin{itemize}
\raggedright
  \item \emph{Say it:} \emph{la turbolenza}
  \item \emph{Say it:} \emph{la turbolenza}, once more
  \item \emph{Say it:} a sentence of your own with \emph{la turbolenza}
  \item \emph{Recall:} say the Italian for to take off, then the Italian for to land, then say \emph{la turbolenza} again
\end{itemize}

\begin{tcolorbox}[breakable,colback=teal!4,colframe=teal!35!black,title={Before you move on}]
What does la turbolenza mean? (\textquotedblleft{}The turbulence\textquotedblright{}.) Say it once more.
\end{tcolorbox}
\section[allacciare]{allacciare — to fasten, to buckle}

\label{lesson:IT-C428-allacciare}



\begin{quote}
Before the new one: say the Italian for to land, then the Italian for the turbulence.
\end{quote}

\subsection*{You'll want to know: allacciare}
\textbf{allacciare} — \textquotedblleft{}to fasten, to buckle\textquotedblright{}.

\begin{quote}
Allacciate le cinture, per favore.
\end{quote}

\begin{quote}
\emph{Fasten your seatbelts, please.}
\end{quote}

\emph{Slacciare} — to unfasten; \emph{allacciare le scarpe} — to tie your shoes.

\begin{cousinweb}
\emph{allacciare} is built from \emph{a-} and \emph{laccio} 'lace, noose', which is inherited from Latin \emph{laqueus} 'noose, snare' (Lewis \& Short).
\end{cousinweb}

\begin{grammarlens}[title={What you've built}]
One more word for telling, explaining and describing at B1.
\end{grammarlens}

\subsection*{Your turn}
\begin{itemize}
\raggedright
  \item \emph{Say it:} \emph{allacciare}
  \item \emph{Say it:} \emph{allacciare}, once more
  \item \emph{Say it:} a sentence of your own with \emph{allacciare}
  \item \emph{Recall:} say the Italian for to land, then the Italian for the turbulence, then say \emph{allacciare} again
\end{itemize}

\begin{tcolorbox}[breakable,colback=teal!4,colframe=teal!35!black,title={Before you move on}]
What does allacciare mean? (\textquotedblleft{}To fasten, to buckle\textquotedblright{}.) Say it once more.
\end{tcolorbox}
\section[il finestrino]{il finestrino — the window (of a car, train or plane)}

\label{lesson:IT-C428-finestrino}



\begin{quote}
Before the new one: say the Italian for the turbulence, then the Italian for to fasten.
\end{quote}

\subsection*{You'll want to know: il finestrino}
\textbf{il finestrino} — \textquotedblleft{}the window (of a car, train or plane)\textquotedblright{}.

\begin{quote}
Preferisco il posto vicino al finestrino.
\end{quote}

\begin{quote}
\emph{I prefer the window seat.}
\end{quote}

At home it is \emph{la finestra}.

\begin{cousinweb}
\emph{finestrino} is a little \emph{finestra} 'window', which is inherited from Latin \emph{fenestra} (EtymDB).
\end{cousinweb}

\begin{grammarlens}[title={What you've built}]
One more word for telling, explaining and describing at B1.
\end{grammarlens}

\subsection*{Your turn}
\begin{itemize}
\raggedright
  \item \emph{Say it:} \emph{il finestrino}
  \item \emph{Say it:} \emph{il finestrino}, once more
  \item \emph{Say it:} a sentence of your own with \emph{il finestrino}
  \item \emph{Recall:} say the Italian for the turbulence, then the Italian for to fasten, then say \emph{il finestrino} again
\end{itemize}

\begin{tcolorbox}[breakable,colback=teal!4,colframe=teal!35!black,title={Before you move on}]
What does il finestrino mean? (\textquotedblleft{}The window (of a car, train or plane)\textquotedblright{}.) Say it once more.
\end{tcolorbox}
\section[(dialogue)]{Ripasso dei capitoli 421–428}

\label{lesson:IT-R428-ripasso-421-428}



\begin{quote}
Say the Italian for the pier and the landing stage.
\end{quote}

\subsection*{Your turn}
Say these aloud:

\begin{itemize}
\raggedright
  \item the pier, the landing stage, to board, to disembark, to set sail
  \item the boat, the rubber dinghy, the oar, to row, the canoe
  \item the cruise, the cabin, the sailor, the bow, the stern
  \item the buoy, to come ashore, to moor, the lifebelt, the nausea
  \item the terminal, the stopover, the boarding, the shuttle, the time zone
  \item the wheeled suitcase, the hold, the belt, the scales, to pack
  \item the customs officer, to declare, to search, the search, the pass
  \item to take off, to land, the turbulence, to fasten, the window
\end{itemize}

\begin{tcolorbox}[breakable,colback=teal!4,colframe=teal!35!black,title={Before you move on}]
The pier? (\textbf{il molo}.) The landing stage? (\textbf{il pontile}.) To board? (\textbf{imbarcarsi}.) To disembark? (\textbf{sbarcare}.) To set sail? (\textbf{salpare}.) The boat? (\textbf{il battello}.) The rubber dinghy? (\textbf{il gommone}.) The oar? (\textbf{il remo}.) To row? (\textbf{remare}.) The canoe? (\textbf{la canoa}.) The cruise? (\textbf{la crociera}.) The cabin? (\textbf{la cabina}.) The sailor? (\textbf{il marinaio}.) The bow? (\textbf{la prua}.) The stern? (\textbf{la poppa}.) The buoy? (\textbf{la boa}.) To come ashore? (\textbf{approdare}.) To moor? (\textbf{ormeggiare}.) The lifebelt? (\textbf{il salvagente}.) The nausea? (\textbf{la nausea}.) The terminal? (\textbf{il terminal}.) The stopover? (\textbf{lo scalo}.) The boarding? (\textbf{l'imbarco}.) The shuttle? (\textbf{la navetta}.) The time zone? (\textbf{il fuso orario}.) The wheeled suitcase? (\textbf{il trolley}.) The hold? (\textbf{la stiva}.) The belt? (\textbf{il nastro}.) The scales? (\textbf{la bilancia}.) To pack? (\textbf{imballare}.) The customs officer? (\textbf{il doganiere}.) To declare? (\textbf{dichiarare}.) To search? (\textbf{perquisire}.) The search? (\textbf{la perquisizione}.) The pass? (\textbf{il valico}.) To take off? (\textbf{decollare}.) To land? (\textbf{atterrare}.) The turbulence? (\textbf{la turbolenza}.) To fasten? (\textbf{allacciare}.) The window? (\textbf{il finestrino}.)
\end{tcolorbox}
